\section{Active roots and their support}\label{sec:activity}

The product rule turns the plaque measures into a directed graph at each
augmented state.  We first show that its arrows compose and describe the
support of the positive-root plaques.  We then realize this graph by
measurable indicators on the original spaces, retaining the long-wall
comparison needed for the entropy argument.

Write $\mathsf P$ for the law of the limiting augmented state
$\omega=(x,(P^B)_B)$.  For $j\ne k$, define
\[
 a_{jk}(\omega)
   =\mathbf 1_{\{m^{U_{jk}}_\omega\text{ is not supported on }\{1\}\}}.
\]
We call $jk$ \emph{active} when $a_{jk}=1$.  This definition does not depend
on the representative of the projective measure.  It is measurable:
after normalizing by a positive integral against a bump near the identity,
activity is detected by a countable family of nonnegative continuous
tests supported away from the identity.  All assertions below are made
on a common set of full $\mathsf P$-measure on which the comparison and
product rules hold.

\begin{proposition}[Composition and positive-root support]
\label{prop:activity}
Almost surely, activity is transitively closed:
\[
 a_{pq}(\omega)a_{qr}(\omega)
       \bigl(1-a_{pr}(\omega)\bigr)=0
       \qquad(p,q,r\text{ distinct}).
\]
Let $b\in L$ be regular, let $V$ be the pattern group of all
$b$-positive roots, and let $V_\omega^+$ be the subgroup generated by
the active $b$-positive root groups.  These active roots form a pattern,
and
\[
                   \supp m^V_\omega\subset V_\omega^+.
\]
Both conclusions hold simultaneously for every regular $b$.
\end{proposition}

\begin{proof}
First observe that Theorem~\ref{thm:product} factors the measure of any
pattern group into its one-dimensional root measures, in a suitable
order.  Indeed, the cone generated by its roots is pointed.  Choose a
root $\alpha$ on an exposed extremal ray.  No other root in the pattern
is proportional to $\alpha$, so a diagonal logarithm can vanish on
$\alpha$ and be strictly negative on every remaining root.  Removing
$\alpha$ leaves a normal pattern subgroup $W$: a sum of two roots in
the original pattern cannot equal $\alpha$, since extremality would
force both summands onto its ray.  Thus the hypotheses of
Theorem~\ref{thm:product} hold for $V=WU_\alpha$.  Iterate, allowing
$W=\{1\}$ at the last step.

To prove transitivity, fix distinct $p,q,r$ and consider the
three-root Heisenberg group generated by $U_{pq},U_{qr},U_{pr}$.
Suppose that $a_{pr}=0$.  Factoring first with $U_{pq}$ as the final
factor, and then with $U_{qr}$ as the final factor, gives two
descriptions of the same projective measure.  Write
$u_{jk}(t)=1+tE_{jk}$.  Since the $U_{pr}$ factor is a point mass, the
first description is supported on the matrices
\[
 u_{qr}(t)u_{pq}(s)=1+sE_{pq}+tE_{qr},
\]
whereas the other is the pushforward of the product of the two root
measures under
\[
 u_{pq}(s)u_{qr}(t)
     =1+sE_{pq}+tE_{qr}+stE_{pr}.
\]
If both roots were active, each root measure would give positive mass
to some compact parameter set disjoint from zero.  Their product would
then give positive mass to matrices with nonzero $pr$ entry, contrary
to the first description.  The projective identities apply on compact
ranges containing these sets, so all masses used here are finite.
This proves transitivity.

For the support assertion, use the successive factorization already
proved.  Every inactive root factor is a point mass at the identity;
every other factor is supported on its own active root group.  The
product is therefore supported on the subgroup generated by the active
roots.  Transitivity says exactly that the active positive roots are
closed under root addition, so this subgroup is a pattern group and is
closed.  There are only finitely many positive-root patterns, which
also gives the simultaneous assertion.
\end{proof}

The entropy argument will use plaques enlarged by the diagonal group.
The following consequence retains the activity at the sampled state.

\begin{lemma}[Support on plaques enlarged by $A$]\label{lem:va-support}
With $V$ and $V_\omega^+$ as in Proposition~\ref{prop:activity}, the
limiting $VA$-plaque probability, written in displacement coordinates
from its sampled point, is supported on $V_\omega^+A$ within its chart.
\end{lemma}

\begin{proof}
Use a product box with coordinates $\psi(q)va$, with product coordinate
domains.  Local $A$-invariance gives, for each $i$, a product
decomposition
\[
                d\mu_i(q,v,a)=d\nu_i(q,v)\,da
\]
inside the box, where $da$ is Haar measure restricted to its coordinate
box and the normalization is absorbed into $\nu_i$.  For $VA$ plaques
the transverse coordinate is $q$; for $V$ plaques it is $(q,a)$.
Choose a common $q$ grid and supplement it by an $a$ grid for the latter
conditioning.  Conditional on a $q$ cell, the $VA$ probability is the
product of a probability on $v$ and normalized Haar measure on $a$.
Conditioning also on an $a$ cell leaves the $v$ probability unchanged.
These are exact identities before taking a limit.  The interior-depth
comparisons permit these compatible grids in the defining data, and
the identities pass to the joint limit by continuous tests and the
null-boundary convention.

At a sample $x=\psi(q_s)v_sa_s$, the $V$ displacements corresponding to
varying $v$ at fixed $a_s$ are
\[
                         a_s^{-1}v_s^{-1}va_s.
\]
By Proposition~\ref{prop:activity}, their probability is supported on
$V_\omega^+$.  A general $VA$ displacement is
\[
 a_s^{-1}v_s^{-1}va
   =\bigl(a_s^{-1}v_s^{-1}va_s\bigr)\bigl(a_s^{-1}a\bigr),
\]
and hence belongs to $V_\omega^+A$ almost surely.  The argument applies
on inner chart boxes; exhausting them gives the statement on the whole
plaque box, whose boundary has zero probability.
\end{proof}

\begin{lemma}[Indicators along the sequence]\label{lem:activity-indicators}
Fix an interior scale $R_i$ allowed in the long-wall comparison of
Proposition~\ref{prop:plaque-comparison}, with ratio-separated depths on
both sides of it inside $(l_{-,i},l_{+,i})$.  There exist measurable
functions $g_{jk,i}:X\to\{0,1\}$ such that
\[
 \operatorname{Law}_{\mu_i}\bigl(Z_i(x),(g_{jk,i}(x))_{j\ne k}\bigr)
 \ \Longrightarrow\
 \operatorname{Law}_{\mathsf P}
       \bigl(\omega,(a_{jk}(\omega))_{j\ne k}\bigr).
\]
They satisfy asymptotic transitivity:
\[
 \int_X g_{pq,i}(x)g_{qr,i}(x)
                  \bigl(1-g_{pr,i}(x)\bigr)\,d\mu_i(x)=o(1)
       \qquad(p,q,r\text{ distinct}),
\]
and, for every fixed $C<\infty$,
\begin{equation}\label{eq:wall-invariance}
 \sup_{\substack{b\in\ker\beta_{jk}\\ \|b\|\le CR_i}}
 \int_X\bigl|g_{jk,i}(xa(b))-g_{jk,i}(x)\bigr|\,d\mu_i(x)
                         \longrightarrow0.
\end{equation}
\end{lemma}

\begin{proof}
Let $\mathcal S$ be the separable metric state space used for $Z_i$,
and put $A_{jk}=\{\omega:a_{jk}(\omega)=1\}$.  Choose
$\epsilon_r\downarrow0$.  For each root and each $r$, there is a
$\mathsf P$-continuity set $D_{jk,r}\subset\mathcal S$ such that
\[
              \mathsf P(D_{jk,r}\mathbin\triangle A_{jk})
                      <\epsilon_r.
\]
For completeness, use a countable family of continuous coordinate
tests generating the Borel sets of $\mathcal S$.  For each test choose
a countable dense family of thresholds outside the atoms of its
distribution under $\mathsf P$.
The finite Boolean combinations of the resulting threshold sets have
null boundary and approximate every measurable set in
$\mathsf P$-measure.  Thus the approximations may be made using
finitely many data tests at each stage.

For fixed $r$, set
$g^{(r)}_{jk,i}(x)=\mathbf 1_{D_{jk,r}}(Z_i(x))$.
Weak convergence of the states and the null boundaries give joint
convergence of the states and these finitely many indicators.  Since
the exact indicators are transitive, the limiting transitivity defect
for any triple is at most $3\epsilon_r$.

We next prove the uniform wall estimate at this fixed approximation
level.  Fix a root, $C$, and any sequence
$b_i\in\ker\beta_{jk}$ with $\|b_i\|\le CR_i$.  Both marginals of
\[
                    \bigl(Z_i(x),Z_i(xa(b_i))\bigr)
\]
have the ordinary state law, because $\mu_i$ is $A$-invariant.
Consequently these paired laws are tight.  In every subsequential
limit their two marginals are $\mathsf P$, and
Proposition~\ref{prop:plaque-comparison} identifies the two projective
$U_{jk}$ measures.  Their exact activity indicators are therefore
equal.  Applying the continuity-set approximation separately in the
two marginals shows that every such limit has
\[
 \mathbb E\bigl|
       \mathbf 1_{D_{jk,r}}(\omega)
        -\mathbf 1_{D_{jk,r}}(\omega')\bigr|
                    \le 2\epsilon_r.
\]
The boundary of the tested event is contained in the union of the two
marginal boundaries, so the corresponding prelimit expectations
converge along this subsequence.  Since the sequence $b_i$ was
arbitrary, a sequence chosen within a vanishing error of the supremum
proves
\[
 \limsup_i\ \sup_{\substack{b\in\ker\beta_{jk}\\\|b\|\le CR_i}}
 \int_X\bigl|g^{(r)}_{jk,i}(xa(b))-g^{(r)}_{jk,i}(x)\bigr|\,d\mu_i(x)
                      \le2\epsilon_r.
\]
This argument uses limits of paired states even when $b_i$ diverges.

Finally choose $r=r(i)\to\infty$ sufficiently slowly.  At stage $r$
impose the preceding estimates for all roots, all triples, and all
integers $1\le C\le r$, together with the first $r$ tests of a
convergence-determining family for the joint state and indicator law.
The limiting laws at fixed $r$ differ from the exact activity law by
at most the sum of the finitely many approximation errors.  A
diagonal choice therefore gives the asserted joint convergence,
vanishing transitivity defects, and \eqref{eq:wall-invariance} for
$g_{jk,i}=g^{(r(i))}_{jk,i}$.  Bounds for integral $C$ imply the bounds
for every fixed finite $C$.
\end{proof}

We conclude with the relation between activity and pieces having
countable intersections with root plaques.
It is this relation that will exclude proper homogeneous components
once all roots have been shown to be active.

\begin{lemma}[Inactivity on pieces with countable plaque sections]\label{lem:transverse-inactive}
Let $S\subset X$ be a Borel set whose intersection with each local
$U_{jk}$ plaque is countable.  Then
\[
            a_{jk}(\omega)=0
         \quad\text{for $\mathsf P$-almost every state with }x\in S.
\]
In particular, this applies to the fixed-point pieces of
Lemma~\ref{lem:fixed-pieces} for roots crossing distinct eigenvalues:
each such piece meets a sufficiently small root plaque in at most one
point.  It also applies to countable unions of these pieces.
\end{lemma}

\begin{proof}
Work first in a $U_{jk}$ flow box with coordinates $(q,h)$ and plaque
probability $P^B$.  For each $q$, the section
\[
                  S_q=\{h:\psi(q)h\in S\}
\]
is countable.  The sampling identity for this single chart states
that the conditional law of the sampled $h$, given $(q,P^B)$, is
$P^B$.  A probability measure assigns mass zero to those points of a
countable set which are not its atoms.  Consequently
\[
 \mathsf P\{x\in S\cap B,\ P^B(\{h\})=0\}=0.
\]
No conditioning on the other plaque probabilities is used here.
Projective agreement in the chart identifies an atom at the sampled
$h$ with an atom at the identity of $m^{U_{jk}}_\omega$.  A countable
flow-box cover therefore shows that this root measure has an atom at
the identity almost surely over $S$.

It remains to show that, under the full state law, an atom at the
identity forces the entire root measure to be supported there.
Identify $U_{jk}$ with $\R$ by $t\mapsto1+tE_{jk}$.  Choose a fixed
$b\in L$ with $\lambda=e^{\beta_{jk}(b)}>1$.  Fixed-motion equivariance
in Proposition~\ref{prop:plaque-comparison}, together with
$A$-invariance of the prelimit sampling law, says that the projective
law of $m^{U_{jk}}_\omega$ is invariant under the pushforward by
$t\mapsto\lambda t$.

The event that this measure has an atom at zero is itself invariant
under dilation.  If that event has positive probability, condition
the projective-measure law on it and normalize each measure $m$ by
its atom: $\nu=m/m(\{0\})$.  This is a measurable normalization of a
locally finite measure, with $\nu(\{0\})=1$, and its law remains
dilation-invariant.  For every compact interval $K$ disjoint from
zero,
\[
 (t\mapsto\lambda^n t)_*\nu(K)
         =\nu(\lambda^{-n}K)\longrightarrow0.
\]
Indeed, $\lambda^{-n}K$ lies in a shrinking punctured neighborhood of
zero, whose mass tends to zero by local finiteness.  Dilation
invariance makes each quantity in this display have the same
distribution as $\nu(K)$.  Almost-sure convergence to zero therefore
forces $\nu(K)=0$ almost surely.  A countable family of such intervals
covers $\R\setminus\{0\}$, and hence $\nu=\delta_0$ almost surely.
The atom event of probability zero is harmless.  Combining this
conclusion with the first part proves the lemma.
\end{proof}
